\section{Matrix Exponentials and Twists}\label{sec:pl:exp-twists}

So far a pose $T\in\SEtwo$ has been a static object: where the robot \emph{is}. A controller, an odometry filter or a motion planner needs the dynamic counterpart: a description of how the pose \emph{changes} that can be scaled, integrated, compared and inverted. The naive choice, ``add a small $(\Delta x,\Delta y,\Delta\theta)$ to the pose'', fails for the same reason it fails in 3-D: poses do not live in a vector space, and a robot that turns while it drives does not move along a straight line. This section builds the right object from one physical picture: a robot repeating the same tiny body-frame step over and over.

\begin{keyidea}
A constant body-frame motion is generated by one matrix $[S]\in\setwo$, the \emph{twist}: it holds the forward velocity and the turn rate per unit of ``motion amount'' $\lambda$. Repeating an infinitesimal step $I+\delta\lambda\,[S]$ infinitely often gives the exact finite motion $T(\lambda)=\exp(\lambda[S])$. In the plane this is always a rotation about a single point, the turn centre $C=Jv/\omega$, or a pure translation. The logarithm goes back from a pose to the twist, but only up to a choice of branch. It recovers the \emph{shortest} generator, not the motion history.
\end{keyidea}

\input{\subfix{sections/05-exponentials-twists/01-small-step}}
\input{\subfix{sections/05-exponentials-twists/02-generator}}
\input{\subfix{sections/05-exponentials-twists/03-exponential-limit}}
\input{\subfix{sections/05-exponentials-twists/04-planar-rotation}}
\input{\subfix{sections/05-exponentials-twists/05-turn-centre}}
\input{\subfix{sections/05-exponentials-twists/06-twists}}
\input{\subfix{sections/05-exponentials-twists/07-logarithms}}
\input{\subfix{sections/05-exponentials-twists/08-lie-bracket}}
\input{\subfix{sections/05-exponentials-twists/09-in-practice}}
